\section{Dining with Dung Beetles}

The diversity of opinions on fundamental issues has always perplexed and intrigued me. In Euclidean geometry, no one would dispute the Pythagorean theorem, at least once he understood it. No scientist would deny that water, steam and ice are the same chemical compound, provided he understands chemistry. But when it comes to political and religious issues, opinions abound. Moreover, people feel compelled to argue for their respective positions in the usually vain attempt to ``convert'' another to that point of view. Yet, that is, in most cases, more difficult than convincing a dung beetle to change his diet.

I am not talking about matters of taste, which few would argue about, but about fundamental issues. It is curious that the average person never holds an opinion about a theorem of algebraic topology or a finding in physical chemistry. Those sciences are clearly beyond his competence and he knows it. Nevertheless, although politics and metaphysics are more difficult, the average person somehow feels entitled to strongly held views on those topics. Somehow, and I always see this in the sophists that offer their opinions on TV programs, it is sufficient to being a sentence with ``I think'', without ever considering why or how they think.

Other factors are ignorance of facts, the use of logical fallacies or mistakes, and oversimplifications that would reduce movements that have lasted for many centuries among various peoples to a simple slogan. Even taking those factors into account, there is still something left over, something inexplicable on the assumption that there is a common human nature.

\paragraph{Evola, Tradition, Spiritual Races}


Someone sent me a video of an interview by a ``professor of tradition'', if you can imagine such a being. The interviewer was a woman who giggled throughout the interview. As expected, the professor mentioned \textbf{Rene Guenon}, \textbf{Ananda Coomaraswamy}, and \textbf{Frithjof Schuon} as the founders of the ``school'', neglecting any mention of \textbf{Julius Evola}. As we have demonstrated here, Guenon and Coomaraswmay respected Evola's work, which cannot be neglected.

Nevertheless, we see some mistakes he made which have unfortunately derailed discussions among men of tradition. There is his idea of ``regality'' apart from priests or warriors. Coomaraswamy demonstrated Evola was incorrect about that due to a misreading of ancient Vedic texts. This ties in to his anti-clerical attitude, which he extends to a diminution of the priestly caste. Had he relied more on Vico than Bachofen, he may have gotten his early prehistory more exactly. The most serious error was equating action with contemplation, which cannot be done consistently. Most recently we saw that in the essay on Plotinus, who was certainly a contemplative rather than a man of action.

His idea of spiritual races, on the other hand, has been insufficiently explored. This may be due to its reduction to biological race, which is not the case at all. Most people, it seems to me, try to make a point based on the notion that an ``idea'' has causative power. For example, in some groups, you will hear how ``Christianity'', for example, destroyed the West or something similar. Ideas have no being except as they manifest in a mind. Hence, ideas can cause nothing; only men holding an idea can be causative provided they act on it.

Ha, that just brings us back to the question that opened this article. Why does a man believe one thing rather than another, and with such stubbornness? Evola says that men are not all alike in their interiority, but rather they possess a ``spiritual race'', which are distinct from one another. In the theory of cosmic cycles, time is not uniform, but it has qualities that differ along the cycle. Hence, by the law of affinity, different spirits will be attracted to the point in the cycle which is compatible with their being. Guenon referred to that as being ``compossible'', i.e., the simultaneous manifestation of possibilities can occur only when they are mutually compatible (not necessarily harmonious).

So in our time, a particular race is manifesting more and more. This is the demeterian race (after the goddess Demeter); member of this race will have a particular worldview that is very hard to shake since it is so fundamental to their being. This is not to say that more atavistic races do not exist, such as the solar-olympian race. These will feel out of place, however. There can be no compromise between the two. The latter is anchored to a Traditional view, while the former is intent on eradicating every trace of it.

\paragraph{Matriarchy}


Relying on \textbf{Otto Weininger} and especially \textbf{J J Bachofen}, Evola brings up the ideas of the masculine and feminine much more than the other early writers on Tradition. He masterfully wove these concepts with the more traditional ideas of castes and cycles in his \emph{Revolt}. Clearly, Gornahoor is not endorsing it pure and simple, because of the errors already mentioned. Nevertheless, the regression of castes is related to the feminization of the higher caste as it reverses its role with the one that had been beneath it. The movement to matriarchy is in full swing at this point as the demeterian race is firmly in control and the higher spiritual values of the solar races are mocked, demeaned, or prohibited.

Matriarchy is marked by the urge to control: what you can eat and drink (e.g., Mrs. Obama with school lunches, Bloomberg in New York over sweet drinks); the desire to prohibit the bearing of arms which is always the prerogative of the free man; the universalization of health care, which will become a prime tool for control; the self-monitoring of free speech. Only the expression of error may be prohibited, if necessary, and truth is a defense. In the engulfing matriarchy, truth is no defense, and the various prohibitions are arbitrary.

In addition to control, the other mark of matriarchy is the pursuit of pleasure. Hence, there is the trend to promiscuity, miscegenation, and eventually polyamory. No sexual expression should be taboo and the deliberate choice of sacrifice in pursuit of something higher is considered absurd when a pleasurable choice is available.

In his important book Eros and Civilization \textbf{Herbert Marcuse} laid out the path to this world perfectly. Obviously, the first step is to overthrow patriarchy which is, in his view, the cause of repression. Moreover, it caused surplus suppression above and beyond that need to maintain a stable society. The guilt associated with that, as represented by Oedipus, must be overcome. Through advances in technology, there will be less need for work, thus freeing up the body for pleasure. It will become polymorphous perverse and will be able to experience sexual pleasure in ways that the patriarchy had discouraged or prohibited. Marcuse expounded the perfect philosophy for the demeterian race.


\flrightit{Posted on 2013-03-20 by Cologero}

\begin{center}* * *\end{center}

\begin{footnotesize}
\begin{sffamily}
\texttt{Avery Morrow on 2013-03-21 at 03:58 said:}

At a certain American university, all students are asked to read an essay that opens with these words:

\emph{There is a pathos in television dialogue: the rapid exchange of monologues that fail to find the issue, like ships passing in the night; the reiterated preface, ``I think that . . .,'' as if it mattered who held which opinion rather than which opinion is worth holding; the impressive personal vanity that prevents each ``discussant'' from really listening to another speaker and that compels him to use this God-given pause to compose his own next monologue; the further vanity, or instinctive caution, that leads him to choose very long words, whose true meaning he has never grasped, rather than short words that he understands but that would leave the emptiness of his point of view naked and exposed to a mass public.}

You can't fault people for imitating the style of rhetoric their TVs have taught them for 50 years.

\hfill

\texttt{Radbod on 2013-03-21 at 05:01 said:}

Wonderful, inspiring article. I may, however, utter my concern regarding the statement that ``ideas have no being except as they manifest in a mind'' and hence ``can cause nothing''. For the authors cited above ideas had to be seen as belonging to the realm of a higher reality. If it is accurate to say that only men holding an idea can be causative provided they act on it, might then be a question of the viewpoint.

\hfill

\texttt{Mihai on 2013-03-21 at 07:14 said:}

1. I find it highly ironic that some Italians claim that ``Evola is Marcuse... only better''. This actually shows how little his work is understood and that few ever go any deeper then the opposition to both capitalism and marxism (which was also endorsed by Marcuse, perhaps the only point in common).

2. What he said and what you summerize here regarding the ``demetarian race'' is highly interesting. I recently spoke with someone with those features and I can say it was really hard getting through in any way to him. He couldn't even conceive any benefit from not giving in to pleasurable impulses and why would anyone reject hedonism, being convinced that it is the most desirable way of life.

The question is this: is such a ``demeterian'' race something fixed, unchangeable, irreversible ?

The answer to this should be ``no''. Since it is a degeneration, it has no independent existence, which means that the process can always be reversed, the polarities brought back to normal, even if in a small degree.

\hfill

\texttt{Cologero on 2013-03-21 at 08:06 said:}

Good observation, Radbod, but I am actually assuming the viewpoint you mention. Perhaps I was unclear. Yes, ideas come from a higher reality, which has many degrees, hierarchically arranged. As spirit, however, these ideas have no compulsion since spirit is free. The human race is flooded with ideas of various quality, some higher, some lower which we then assume are our ``own'' ideas. The ideas are in conflict, hence there is an inner spiritual combat, at least to the extent we can detach from self-identifying with them. But not even the devil can control our will, he works by suggestion. The point is that certain constellations of ideas appeal to different kinds of people, which we could rightly call their ``spiritual race''.

\hfill

\texttt{J. Alexander Maximilian on 2013-03-21 at 09:08 said:}

So then the ``Demetrian'' race would be more prone to technological/ external ``progress'' as opposed to organic/ interior (spiritual) progress.

This is interesting as my cousin married a native American woman who is intent on preserving their traditions while at the same time promoting with much zeal the progressive initiatives of American liberalism, which seems to arise from her dislike for traditional Western culture.

She even told me once that my kind was dying and the feminine cycle was ``returning''. It seems to me that much of the ``demetrian'' spirit itself is a creation of their own emotions such as contempt, and feeling of entitlement.

\hfill

\texttt{Radbod on 2013-03-21 at 09:18 said:}

Thank you for your expedient reply. I fully agree.

Please allow me to take this opportunity to express my high esteem for your endeavours. I do not only highly appreciate your translations of previously untranslated works of Evola, since my command of the Italian language unfortunately still leaves a lot to be desired, but also very much your own comments.

\hfill

\texttt{Jason-Adam on 2013-03-21 at 11:18 said:}

I am an admirer of Julius Evola, I find his writings to be very enjoying and full of insights, he is a Traditionalist I say because he does believe in a universal primordial single truth. However, what troubles me about him is something far beyond what ever errors any fallible human can make -- it is what he told Sheikh Pallavicini in 1951 that he was more interested in politics than in spirituality. Too many followers of Evola have that same attitude and when I first came across Traditionalism, I thought it was merely a political ideology like communism or fascism. Even many Traditionalists can not speak of anything but politics and have no clue as to the transcendent.

\hfill

\texttt{Jason-Adam on 2013-03-21 at 11:23 said:}

Now as to the masculine-feminine solar-lunar dialectic, it is not as clear cut as Evola thought. Amaterasu no kami, the progenitor of the Imperial family of Japan, is a sun goddess. Evola wrote once it was a confusion when the mediaeval church, the mother church, called itself the sun yet the Japanese would have had no objections. I am not an expert in these things but perhaps there is a solar femininity ?

Also, I am surprised no one brings the point that the crusades were a conflict between the Christians fighting under the banner of the solar cross with the Mohammedans fighting under the crescent of Allah the moon goddess -- the Christian doctrine of free will which reads as much more masculine and titanic than the passive fatalism of Mohammedanism underscores the comparison.

\hfill

\texttt{Cologero on 2013-03-21 at 13:02 said:}

It should not be troubling, per se, if his mentality was more like the kshatriya caste. In that case, the political, or the world of action, would be his primary area of interest. Evola, unfortunately, goes a step too far when he tries to portray ``action'' and ``contemplation'' as parallel ways of equal merit. First of all, there is the obvious point that one should not act unless one knows what he is doing, and for what purpose. Beyond that, there is a rupture in the hierarchical relationship, with contemplation superceding action: ``parallel'' and ``equal'' have no place in that hierarchy.

I think it is of value for the kshatriya type to write about Tradition; that is why Evola is more accessible to more people than a Guenon could be. However, he cannot reverse the true order of things the way he often does. Fundamentally, Evola, judging from his comments, does not quite understand what ``contemplation'' really is. Often he confuses it with discursive reasoning of the professorial type; clearly that is false and the essay on Plotinus proves it, despite the attempt to ``nietzscheanize'' him as you noticed.

The irony is that the political is feminine in relationship to the metaphysical.

\hfill

\texttt{Jason-Adam on 2013-03-21 at 16:21 said:}

Evola does fill in a gap in that Guenon wrote totally from a priestly perspective, a perspective that Schuon and others followed. What saddens me though is that there is a divide between groups that are solely metaphysical and groups that are solely political when if there is going to be a rectification of the world, the metaphysical groups need to command and direct the political soldiers to fight the holy war. De Giorgio expressed the proper relationship between the castes in the best way possible.

Can you imagine what would have happened if say all the neo-Evolian ``terrorist'' sects of the 1970s had had proper spiritual leadership who could have used the devoted militants in a wise manner ? We today might be in a better world... .

Isn't the irony of ironies that German idealism and vitalism which worships the masculine is at its core feminine ? I wonder if the homosexuality associated sometimes with this train of thought (to his credit while Evola was not pro-homosexual but some of his Mannerbund writings have been used by gays to support their agenda) is a sign of its innate feminine spirit ?

\hfill

\texttt{Jason-Adam on 2013-03-21 at 16:24 said:}

This is tangental but what caste is a Templar ? they seem to combine features of both the priestly and warrior natures.

Also, what caste does Dugin belong to ? In all the writings of his I've read he seems to have a poor grasp of the nature of caste.

\hfill

\texttt{J. Alexander Maximilian on 2013-03-21 at 17:17 said:}

Templars and other Crusading orders were composed of warrior/ priests, monks. This can also be found in eastern religions like Buddhism. As for caste it would seem they exist in a between place. In their renunciation of the world, even though they were warriors they were monastics who belonged entirely to the Church, therefore priestly caste would be most appropriate.

Dugin? I cannot elaborate on that as it seems to me he's very politically motivated, not really traditional. That's my POV on Dugin.

\hfill

\texttt{Saladin on 2013-03-21 at 19:56 said:}

``will be able to experience sexual pleasure in ways that the patriarchy had discouraged or prohibited. ''

Not only Patriarchy discouraged and prohibited some of these perversions in an innocent manner, but also it was completely unaware of many more which moderns delight in. Many of the sexual mores of today would be completely unknown to normal man of Tradition (even in his wildest of imaginations).

\hfill

\texttt{Saladin on 2013-03-21 at 20:12 said:}

Good point Mihai! Every man has ``demeterian'', chthonic, subterranean, or lunar tendencies to a more or lesser degree. However, the Solar type (or race if you will) does not succumb to it. This resistance to lower tendencies is the meaning and purpose of the Greater War. The Lunar type, on the contrary, rejoices in losing himself in the mire of promiscuity. The Lunar man can never reveres his tendencies and as you say bring things ``back to normal''. This overcoming is reserved only for the virile Uranic man.

\hfill

\texttt{Saladin on 2013-03-21 at 20:56 said:}

Jason-Adam, you criticize Evola for being too simplistic to designate a ``solar-lunar dialectic'', yet you yourself fall into the same dualistic mindset. I hope that my being a moslem doesn't make me too biased in saying you are wrong to ascribe to the moslems (and they are not ``Mohammedans'' as you suggest since Islam unlike christianity does not revolve around the person of Mohammad) a lunar character. The early Muslim community did not really have a symbol. During the time of the Prophet Muhammad (peace be upon him), Islamic armies and caravans flew simple solid-colored flags (generally black, green, or white) for identification purposes. In later generations, the Muslim leaders continued to use a simple black, white, or green flag with no markings, writing, or symbolism on it. It wasn't until the Ottoman Empire that the crescent moon and star became affiliated with the Muslim world.

Also, I am afraid you show your ignorance (and maybe Prejudice?!) when you talk about ``the Christian doctrine of free will'' and ``the the passive fatalism of Mohammedanism''. You have read too much of Houston Stewart Chamberlain, Alfred Rosenberg, and other Racialists of 19th and 20th century. Both concepts of Free-will and Predestination are present in theologic debates in Islam and Christianity alike.

Regarding the Crusades I agree with Evola that both camps were Traditional and Solar (consult his Revolt against the Modern World). He states in his Metaphysics of War: ``the ancient Arab Knight ascended to the same supra-traditional point which the Crusader Knight attained by his heroic asceticism.''

\hfill

\texttt{Mihai on 2013-03-22 at 04:37 said:}

Well I say that even for the Lunar type there is, or at least was, a possibility of overcoming his promiscuity and attaching himself, at least indirectly, to a higher order.

In Eastern Christianity it is said there are three paths by which one can achieve salvation.

The first (in ascending order) is the path of the slave- the man who follows God for fear of Hell and perpetual punishment.

The second is that of the mercenary or hired man who follows God with the eternal rewards of Heaven in mind.

The third is the path of the son, who follows God because it is the only conceivable life and he understand and participates fully into this truth, being in the Kingdom of Heaven as a son is on his father's estate.

Clearly, these three correspond to the three tendencies of tamas, rajas and sattwa.

Clearly, even for the tamasic person, the inferior man, there is still a hope, regardless how limited it is.

Of course, nowadays that Hell is considered merely a fantasy and no ``respectable intellectual'' would pay any attention to it, even that has been taken away.

\hfill

\texttt{Saladin on 2013-03-22 at 07:43 said:}

Absolutely! God's mercy includes all men, provided, as you mentioned, they place themselves within a regular tradition.

\hfill

\texttt{Jason-Adam on 2013-03-22 at 11:02 said:}

My point was that the solar-lunar dialectic is too simplistic precisely because many allusions can be made that are only based on surfaces. My point re. the cross vs the crescent was meant to show the flaw in the arguement.

I do think there is a meaning to solar and lunar but that the meaning is something deeper than what Bachofen thought. Perhaps someone else can enlighten me.

\hfill

\texttt{Avery Morrow on 2013-03-22 at 11:18 said:}

``... yet the Japanese would have had no objections. I am not an expert in these things but perhaps there is a solar femininity ?''

Read my upcoming book and learn all about how the medieval Japanese believed that Amaterasu was actually male. Evola was right!

\hfill

\texttt{Cologero on 2013-03-22 at 12:55 said:}

First of all, the language is symbolic and is emotionally charged to many. Who wants to be lunar or feminine? Better would be the terms active and passive. The second and more important issue is that the terms are relative and not absolute. E.g., as we pointed out (following Coomaraswamy) the kshatriya is passive in relationship to the brahmana, but active in respect to the lower castes. So it is always ``lunar'' in respect to what ... 

\hfill

\texttt{J. Alexander Maximilian on 2013-03-22 at 13:12 said:}

Excellent explanation. No human is absolutely one or the other. But one would have more influence over an individual, society.

As was pointed out one can be within Solar lineage and have a lunar position as a contemplative mystic.

\hfill

\texttt{J. Alexander Maximilian on 2013-03-22 at 13:14 said:}

Sounds good. One of my favorite times and place.

\hfill

\texttt{Izak on 2013-03-22 at 14:29 said:}

``The irony is that the political is feminine in relationship to the metaphysical.''

Yes, and it is kind of interesting how he attempts to maneuver around that pesky detail in his discussion of John Woofroffe and the nature of sakti in his own The Yoga of Power... although it is a good book nevertheless for those interested in more worldly inner self transformation.

\hfill

\texttt{Cassiodorus on 2013-03-24 at 13:30 said:}

I just came across a rather provocative essay on Guenon that I think the Gornahoor community will find interesting. It was written by the author of ``Inner Christianity'', Richard Smoley, a scholar of Western esotericism. The title is '' Against Blavatsky: Rene Guenon's Critique of Theosophy''.

\textit{theosophical.org/publications/quest-magazine/1696}


\hfill

\texttt{Cologero on 2013-03-24 at 20:27 said:}

How exactly does that link relate to the post in question?

\hfill

\texttt{Avery on 2013-03-24 at 21:43 said:}

I've stumbled on that article more than once. My favorite part is when the author seems to admit that Guénon does have a better interpretation of the East that actually exists.

``Here Guénon stands on firmer ground. The concept of an evolving humanity in an evolving universe is very difficult to find in traditional Eastern texts.''

But in the end he has to side with the fantastical East of Mme. Blavatsky, because castes are just so... old, and backwards!

\hfill

\texttt{Saladin on 2013-03-24 at 22:02 said:}

Thanks for the link Cassiodorus, although the article does not tell us anything which we already did not know.

What is amazing is the author's incapacity to distinguish what the ancients meant by transmigrations of souls and the modern idea of reincarnation. He still insists that Plato's symbolic description of postmortem condition should be taken literally as Theosophists and other New Agers do. In order to truly grasp what Pythagoras, Plato, and others gnostics meant by transmigration one needs to understand the Multiple States of Being.

\hfill

\texttt{Cassiodorus on 2013-03-24 at 23:26 said:}

Cologero,

It doesn't relate. Apologies for the derail.

\hfill

\texttt{Cologero on 2013-03-24 at 23:43 said:}

Cassiodorus, it may have fit in better on a different post. When providing links, it would be useful to point out what in particular is of interest.

\hfill

\texttt{Cavalcare on 2013-04-01 at 13:35 said:}

1. Explain this please ``German idealism and vitalism which worships the masculine is at its core feminine''.

2. Homoeroticism =/= homosexuality.

\hfill

\texttt{Sparrow on 2016-04-15 at 10:26 said:}

``Matriarchy is marked by the urge to control: what you can eat and drink (e.g., Mrs. Obama with school lunches, Bloomberg in New York over sweet drinks); the desire to prohibit the bearing of arms which is always the prerogative of the free man; the universalization of health care, which will become a prime tool for control; the self-monitoring of free speech.''

It's possible that I'm misreading something here, but do you mean to suggest that in a healthy, patriarchal, traditional society, that the disgusting practice of stuffing one's face with poison and becoming obese would be permitted? Or that the low dregs of society would be permitted to bear arms, with which they could threaten the well bred castes, as we've seen in countless revolutions. I'm not sure why you would be appealing to the modern conception of ``freedom,'' the one brought about by the American and French liberal revolutions, with which the decadence of modern society is okay because ``everything goes.''

\hfill

\texttt{Cologero on 2016-04-15 at 23:55 said:}

Wow, Sparrow, right after I've boasted of my proficiency in logic, you come up with this. I'm challenged to unravel your thought process.

Apparently, you are confusing a traditional society with a Puritanical society. For the sake of public order, not every vice would be prohibited by law. For example, the disgusting practice of masturbation would not be illegal. Gluttony is certainly a lesser evil than that. Nevertheless, the Tantras prohibit both the \textit{lewd and the glutton}\footnote{\url{http://www.gornahoor.net/?p=55}} from discipleship.

As for the claim that the right to bear arms is the ``prerogative of the free man'', I did not define the ``free man''. Hence, your objection does not logically follow. In any event, the controlling society would prohibit arms even from the free man. More recently, however, \textit{I defined the ``free man''}\footnote{\url{http://www.gornahoor.net/?p=8535}} as one capable of acting morally.

Keeping in mind that there are various types of people, \textit{One law for the lion and ox is oppression}\footnote{\url{http://www.gornahoor.net/?p=1647}}. The pneumatics would not need a law. The psychics would follow the law, and no law would suffice for the hylics. One must apply force judiciously and only when the force of myths and education fall short.

\hfill

\end{sffamily}
\end{footnotesize}

